\section*{第 \ref{ch:b3:bacteriology} 章——细菌学：生长、生理与遗传}

\begin{solution}{exo:b3:bacteriology:1}
\omterm{def:b3:bacteriology:envelope}{革兰氏阳性}菌：细胞质膜外包一层厚而多层的
\omterm{def:b3:bacteriology:envelope}{肽聚糖}壁，壁中穿插磷壁酸。\omterm{def:b3:bacteriology:envelope}{革兰氏阴性}菌：
\omterm{def:b3:bacteriology:envelope}{肽聚糖}层薄，位于周质中，外为\omterm{def:b3:bacteriology:envelope}{外膜}，
\omterm{def:b3:bacteriology:envelope}{外膜}外层为\omterm{def:b3:bacteriology:envelope}{脂多糖}，并含孔蛋白。结晶紫–碘
复合物被滞留于\omterm{def:b3:bacteriology:envelope}{革兰氏阳性}菌厚而经酒精脱水的细胞壁中，呈紫色；
而在\omterm{def:b3:bacteriology:envelope}{革兰氏阴性}菌中，该复合物随薄壁和受破坏的\omterm{def:b3:bacteriology:envelope}{外膜}被洗脱，
随后染上复染剂的
粉红色。
\end{solution}

\begin{solution}{exo:b3:bacteriology:2}
$10^{6}$ 倍 $= 2^{19.9}$：约 $20$ 代；$T = \qty{8}{h}/20
= \qty{24}{min}$; $\mu = \ln 2/\qty{0.4}{h} = \qty{1.7}{h^{-1}}$.
\end{solution}

\begin{solution}{exo:b3:bacteriology:3}
转化——格里菲思（1928 年）以及埃弗里、麦克劳德和麦卡蒂（1944 年）：
热杀死的有毒力肺炎球菌可使活的无毒力菌株发生转化，
转化因子是 DNA。\omterm{def:b3:bacteriology:hgt}{转导}——津德和莱德伯格（1952 年）：噬菌体 P22
能将\emph{Salmonella}的基因从一株\omterm{def:b3:cancer-biology:hallmarks}{转移}到另一株，
而细菌本身无法通过滤器。\omterm{def:b3:bacteriology:hgt}{接合}——莱德伯格和塔特姆（1946 年）：
两株营养缺陷型\emph{E.~coli}只有在混合培养时，
才能产生原养型重组体，且需要细胞接触，
（戴维斯 U 形管实验）。
\end{solution}

\begin{solution}{exo:b3:bacteriology:4}
青霉素：作用于交联\omterm{def:b3:bacteriology:envelope}{肽聚糖}的转肽酶，
动物细胞不具有该酶。链霉素：作用于细菌核糖体小亚基，
其细菌型与真核型不同。环丙沙星：作用于 DNA 促旋酶，
这是一种细菌拓扑异构酶，人体内不存在。
利福平：作用于细菌 RNA 聚合酶的 $\beta$ 亚基，
真核生物的同类酶不含该亚基。
\end{solution}

\begin{solution}{exo:b3:bacteriology:5}
$S^{*} = K_{S}D/(\mu_{\max} - D)$，$X^{*} = Y(S_{0} - S^{*})$。$D = 0.2$：
$S^{*} = \qty{0.0067}{g/L}$，$X^{*} = \qty{0.80}{g/L}$。$D = 0.6$：
$S^{*} = 0.06$，$X^{*} = 0.78$。$D = 0.79$：$S^{*} = 1.58$，$X^{*} = 0.17$。
发生洗出的条件为 $D = \mu(S_{0}) = 0.8\times 2/2.02 = \qty{0.79}{h^{-1}}$。
\end{solution}

\begin{solution}{exo:b3:bacteriology:6}
$(N/V)_{\text{阈}} = kA_{\text{阈}}/p$。发光器官：$0.01\times
6\times 10^{12}/100 = 6\times 10^{8}$ 个/升 $= 6\times 10^{5}$ 个/毫升。
海水：$10\times 6\times 10^{12}/100 = 6\times 10^{11}$ 个/升
$= 6\times 10^{8}$ 个/毫升——高出一千倍，
海水中从未达到此值。
\end{solution}

\begin{solution}{exo:b3:bacteriology:7}
\qty{25}{mm}：敏感，MIC 低。\qty{12}{mm}：敏感性降低，
属中介类型——按常规剂量可能无效。\qty{0}{mm}：耐药，
细菌一直生长至纸片边缘。药物从纸片向外扩散，其浓度随距离
以可重复的方式下降；细菌停止生长处的药物浓度等于 MIC，
因此抑菌圈半径与 MIC 之间的关系由
已知 MIC 的菌株所建立的标准曲线确定。
\end{solution}

\begin{solution}{exo:b3:bacteriology:8}
该\omterm{def:b3:bacteriology:hgt}{接合}性\omterm{def:b3:genetic-engineering:tools}{质粒}携带一个整合子，其盒式排列中含有
多个耐药基因，另外还带有携带更多基因的转座子；\omterm{def:b3:genetic-engineering:tools}{质粒}
一次\omterm{def:b3:bacteriology:hgt}{接合}即可\omterm{def:b3:cancer-biology:hallmarks}{转移}，并同时表达所有这些基因。
这些基因聚集在一起，是因为对其中任一药物的选择都会维持整个\omterm{def:b3:genetic-engineering:tools}{质粒}，
因此搭便车的基因得以保留（共选择）；整合子捕获盒式结构，
转座子跳入\omterm{def:b3:genetic-engineering:tools}{质粒}；
\omterm{def:b3:cancer-biology:hallmarks}{转移}的单位是\omterm{def:b3:genetic-engineering:tools}{质粒}，而不是单个基因。
\end{solution}

\begin{solution}{exo:b3:bacteriology:9}
单用一种药物：预期耐药细胞数为 $10^{-7}\times 10^{9} = 100$。两种独立药物：
预期数为 $10^{-14}\times 10^{9} = 10^{-5}$。持留菌不是突变体：
一小部分休眠细胞能在任何药物中存活，
并在药物消失后重新生长为敏感种群，因此第二种药物无效；
只有足够长的疗程，在其苏醒时将其杀灭，
或使用对休眠细胞有效的药物，才能奏效。
\end{solution}

\begin{solution}{exo:b3:bacteriology:10}
$X > X^{*}$：消耗速率 $\mu X/Y$ 超过供应，因此 $S$ 降至
$S^{*}$ 以下；此时 $\mu(S) < D$，$\mathrm{d}X/\mathrm{d}t < 0$，
$X$ 回落；随 $X$ 下降，$S$ 又重新上升。两种偏离均被纠正，
该稳态是稳定的。两个物种：每个物种维持存在都需要 $\mu_{i}(S) = D$，
由此确定各自的 $S_{i}^{*}$；营养物质只能取一个值，因此至多一个物种能够存活。
$S^{*}(D)$ 较低的那个物种会将 $S$ 压到另一物种所需水平之下，
另一物种则被洗出；
若两者的莫诺曲线相交，
哪一物种占优可能随 $D$ 而变化。
\end{solution}

\begin{solution}{exo:b3:bacteriology:11}
记密度 $\rho = N/V$。关闭态 $A = p_{0}\rho/k$ 在其值低于
$A_{\text{阈}}$ 时成立，即 $\rho < kA_{\text{阈}}/p_{0}$；开启态 $A = p_{1}\rho/k$ 在其值高于阈值时成立，
即 $\rho
> kA_{\text{阈}}/p_{1}$。当 $kA_{\text{阈}}/p_{1} < \rho <
kA_{\text{阈}}/p_{0}$ 时两种状态都能自洽：已开启的种群因其较高的输出而维持开启，
未开启的则保持关闭。用途：承诺——一旦菌落决定形成生物膜或释放毒素，
密度的短暂下降不会撤销该决定，
且开关是陡峭同步的，
而非渐变的。
\end{solution}

\begin{solution}{exo:b3:bacteriology:12}
(a) 短疗程高剂量：药物水平始终高于突变体的 MIC，
可同时杀灭野生型与一步突变体——有利，但受毒性限制。
(b) 长疗程低剂量：药物水平在数天内停留于选择窗口内，
促使耐药性产生——不利。(c) 症状消失即停药：幸存者是最不敏感的细胞，
药物水平下降时会穿过选择窗口——不利。(d) 联合用药：
细胞必须同时耐受两种药物，
而十亿个细胞中通常不会含有这样的细胞
——有利。
\end{solution}

\begin{solution}{pb:b3:bacteriology:1}
\textbf{1.} $18$ 代：$2^{18} = 2.6\times 10^{5}$ 个细胞，
即每毫升 $2.6\times 10^{3}$ 个。
\textbf{2.} $2\times 10^{9}\times 100 = 2\times 10^{11}$ 个细胞 $=
2^{37.5}$：经过 \qty{12.5}{h}；质量 \qty{0.2}{g}。
\textbf{3.} $6\times 10^{27}/10^{-12} = 6\times 10^{39} = 2^{132}$ 个细胞：
$132$ 代，\qty{44}{h}。营养物质、氧气和废物会使其
在一天之内停止；指数生长总是很短暂。
\textbf{4.} \omterm{def:b3:bacteriology:phases}{稳定期}细胞已拆解核糖体，并抑制了
其代谢酶；分裂之前必须重新合成它们。
指数生长期的细胞则装备齐全。
\textbf{5.} $\mu = \ln 2/\qty{0.333}{h} = \qty{2.08}{h^{-1}}$；当 $T =
\qty{60}{min}$ 时，\qty{0.69}{h^{-1}}；经过 \qty{6}{h} 后 $2^{6} = 64$
而非 $2^{18}$，即少了约 $2^{12} \approx 4000$ 倍。
\textbf{6.} $200/(0.1\times 10^{-6}) = 2\times 10^{9}$ 个活细胞，每
毫升。显微镜还会计入死细胞，以及在
平板上无法生长的细胞；一团聚集体只形成一个菌落。
\textbf{7.} $D = 0.5/1 = \qty{0.5}{h^{-1}}$；$T = \ln 2/0.5 =
\qty{1.39}{h} = \qty{83}{min}$。
\textbf{8.} $S^{*} = 0.01\times 0.5/0.5 = \qty{0.01}{g/L}$；$X^{*} =
0.5\times 1.99 \approx \qty{1}{g/L} = 10^{12}$ 个细胞每升，即每
毫升 $10^{9}$ 个。
\textbf{9.} $0.5\times 10^{12} = 5\times 10^{11}$ 个细胞每小时；$720/1.39
\approx 520$ 代（每月）。
\textbf{10.} $D = 0.95$：$S^{*} = 0.01\times 0.95/0.05 = \qty{0.19}{g/L}$，
$X^{*} = 0.5\times 1.81 = \qty{0.9}{g/L}$。当 $F = \qty{1.1}{L/h}$ 时，$D
> \mu_{\max}$：洗出。
\textbf{11.} 突变体维持的 $S^{*} = 0.01\times 0.5/0.7 = \qty{0.0071}{g/L}$；
在该葡萄糖浓度下，原始菌株的生长速率为 $1.0\times 0.0071/0.0171 =
\qty{0.42}{h^{-1}} < D$，并按 $e^{-0.08t}$ 下降，在
数周内被洗出。突变体随之占据优势。
\textbf{12.} $10^{-9}\times 10^{12} = 1000$ 个有益突变，每
代：适应持续进行，一个新的有利克隆
每隔数百代便席卷种群——\omterm{thm:b3:bacteriology:monod}{恒化器}是一台
演化机器。
\textbf{13.} $kA_{\text{阈}}/p = 0.05\times 6\times 10^{12}/500 =
6\times 10^{8}$ 每升，即每毫升 $= 6\times 10^{5}$ 个细胞。
\textbf{14.} $10^{10}/6\times 10^{5} \approx 1.7\times 10^{4}$ 倍。
\textbf{15.} $20\times 6\times 10^{12}/500 = 2.4\times 10^{11}$ 每升，
$2.4\times 10^{8}$ 每毫升——比海水中的
$10^{2}$ 每毫升高出六个数量级：不发光。
\textbf{16.} 增长十倍需 $3.3$ 次倍增，\qty{6.6}{h}。当 $10^{9}$
每毫升时，种群仍比阈值高出 $1700$ 倍：
发光持续进行。
\textbf{17.} 在海水中无人看到这种光，它也不为任何宿主所用：一个
不发光的突变体节省 \qty{20}{\%} 的代价，并比发光者长得更快。在器官内
乌贼选择发光者——它会驱逐或不喂养不发光的
作弊者——因此这一代价换来了栖身之地。
\textbf{18.} 阻断受体或破坏\omterm{def:b3:bacteriology:quorum}{自诱导物}（\omterm{def:b3:bacteriology:quorum}{群体感应}
淬灭）：细菌仍能存活，但从不启用其毒素，免疫
系统即可清除它们。由于敏感菌与“耐药”菌
生长速度相同——无一被杀死——逃避药物的选择优势
很小，因此耐药性的传播应\omterm{def:b3:bioinformatics:alignment}{比
对}杀菌药物时要慢得多。
\textbf{19.} A：$10^{-8}\times 10^{9} = 10$；B：$100$；两者：$10^{-15}\times
10^{9} = 10^{-6}$。
\textbf{20.} $P_{0}(\text{A}) = e^{-10} = 5\times 10^{-5}$；
$P_{0}(\text{两者}) = e^{-10^{-6}} = 0.999999$。
\textbf{21.} $16 \to 8$：一个半衰期，\qty{4}{h}；$16 \to 1$：四个
半衰期，\qty{16}{h}；在 \qty{4}{h} 至 \qty{16}{h} 的选择窗口内，
每次给药停留 \qty{12}{h}。
\textbf{22.} 每 \qty{12}{h} 给药一次：在第 $4$ 小时至
第 $12$ 小时的窗口内所占时间约为三分之二；十天之内，
A 的一步突变体——开始时共十个——被选出，而
单用 A 则失败。
\textbf{23.} 全程保持高于 \qty{8}{mg/L}：每个半衰期给药一次
（\qty{4}{h}），或每 \qty{12}{h} 给药一次、峰值达 \qty{64}{mg/L}
（$64 \to 8$，经三个半衰期），或持续输注。代价为：
高峰值的毒性，或每天六次给药，而患者
无法坚持。
\textbf{24.} $10^{-4}\times 10^{9} = 10^{5}$ 个持留菌无论
用何种药物都能存活；停药后它们苏醒，并
重新长成完全敏感的种群——即复发。治疗须持续足够长的时间，
以便在其离开休眠时将其杀灭，这正是结核病疗程长达六个月的原因。
\textbf{25.} \omterm{thm:b3:bacteriology:monod}{恒化器}中每毫升 $10^{9}$ 个细胞；发光开启的
阈值为每毫升 $6\times 10^{5}$ 个细胞；$P_{0} = 0.999999$，
即两药疗程中不遇到耐药细胞的概率。
\end{solution}
